% Quantitative Researcher Internship — Two Sigma
\documentclass[10pt,letterpaper]{article}

\usepackage[margin=0.75in]{geometry}
\usepackage[hidelinks]{hyperref}
\usepackage[skip=4pt]{parskip}

\pagestyle{empty}
\pagenumbering{gobble}
\setlength{\parindent}{0pt}

\begin{document}

\textbf{Rojae Mighty}\\
Stony Brook, NY \quad$\cdot$\quad
\href{mailto:rojae.mighty@stonybrook.edu}{rojae.mighty@stonybrook.edu}\\
\href{https://rojaemighty.com}{rojaemighty.com} \quad$\cdot$\quad
\href{https://www.linkedin.com/in/rojaemight/}{linkedin.com/in/rojaemight}

\vspace{8pt}
July 14, 2026

\vspace{8pt}
Hiring Manager\\
Two Sigma\\
New York, NY

\vspace{8pt}
\textbf{Re: Quantitative Researcher Internship}

\vspace{4pt}
Dear Hiring Manager,

I am excited to apply for the Quantitative Researcher Internship at Two Sigma. As an undergraduate physics student at Stony Brook University, I have developed strong quantitative, analytical, and computational skills through research in nuclear and particle physics, machine learning, and financial data analysis. My experience studying geometric scaling and gluon saturation in electron-ion collisions has trained me to investigate complex systems, test hypotheses, and draw reliable conclusions from challenging datasets. I am eager to bring this scientific approach to quantitative research in financial markets.

Through my work at Brookhaven National Laboratory, including the Department of Energy Science Undergraduate Laboratory Internship, I conducted research on nuclear modification in Electron-Ion Collider kinematics using PYTHIA-based Monte Carlo event generation. I utilized an impact-parameter-dependent nuclear modification factor to analyze parton behavior across varying nuclear environments. This work required computational modeling, statistical data analysis, careful evaluation of assumptions, and attention to detail when interpreting uncertain results. My proficiency in C++, Python, Linux, and ROOT has also prepared me to build reproducible analysis workflows and contribute effectively in a technical research environment.

As an undergraduate researcher at Stony Brook's Center for Frontiers in Nuclear Science, I use C++, ROOT, Python, and machine-learning methods to analyze large-scale collision datasets and reconstruct physical quantities. Beyond my academic research, I have developed a market microstructure forecasting project using Python, pandas, NumPy, and scikit-learn to analyze limit order book data. The project includes feature engineering and walk-forward validation, strengthening my understanding of time-series modeling and the importance of evaluating models out of sample. I also completed J.P. Morgan's Quantitative Research program on Forage, gaining experience with derivatives pricing, statistical data analysis, and credit-risk modeling.

Two Sigma's commitment to combining advanced technology, artificial intelligence, and rigorous human inquiry to identify market opportunities strongly resonates with my interest in scientific research. I am particularly excited by the opportunity to collaborate with engineers, quantitative researchers, and data scientists to develop creative, evidence-based solutions to complex economic problems. While my research background is grounded in quantum chromodynamics and accelerator science, the underlying process of modeling complex phenomena, validating results, and communicating technical findings is directly transferable to quantitative finance. I would welcome the opportunity to contribute my curiosity, technical skills, and quantitative problem-solving perspective to Two Sigma.

Thank you for considering my application. I look forward to discussing how my research experience, computational background, and interest in quantitative finance could contribute to the success of your team.

\vspace{4pt}
Sincerely,\\[12pt]
Rojae Mighty

\end{document}
